\subsection{Métodos}

Para el desarrollo de la presente investigación se adoptó una modalidad metodológica de carácter aplicativo, con enfoque cuantitativo y diseño preexperimental. Estas tres características no se trabajaron de forma aislada, sino de manera complementaria a lo largo del proyecto: el componente aplicativo orientó el diagnóstico del flujo de información actual de Bellidental y la construcción de un sistema informático que respondiera a ese diagnóstico, mientras que los componentes cuantitativo y preexperimental permitieron traducir el impacto de dicho sistema en cifras concretas, mediante la medición y comparación del tiempo de gestión administrativa del personal antes y después de su implementación. A continuación se detalla cómo se articula cada una de estas modalidades con el desarrollo específico del proyecto.

\subsubsection{Modalidad de la investigación}

La investigación se estructura bajo componentes metodológicos articulados con las técnicas e instrumentos de recolección de información: enfoque cuantitativo, alcance aplicado, diseño preexperimental, e investigación documental y de campo. Estos elementos orientan el análisis de la problemática identificada en el Centro Odontológico Bellidental y la evaluación de la propuesta desarrollada.

El enfoque cuantitativo permite obtener y analizar datos numéricos relacionados con la carga operativa y los tiempos de gestión. Se emplea principalmente mediante la encuesta dirigida a la secretaria y la ficha de observación cronometrada, cuyos registros permiten comparar los tiempos correspondientes a las actividades administrativas antes y después de la intervención.

El alcance aplicado se orienta a la resolución del problema identificado mediante el desarrollo e implementación de una solución tecnológica. Para ello, la entrevista semiestructurada al personal médico y administrativo permite identificar los procesos, necesidades y dificultades existentes, mientras que la encuesta a los pacientes permite conocer los canales de comunicación utilizados y su disposición hacia el uso de un asistente automatizado.

El diseño preexperimental, de un solo grupo con medición antes y después, permite evaluar los cambios producidos tras la implementación de la solución tecnológica. La ficha de observación cronometrada se aplica en las fases de pretest y postest sobre los mismos procesos: registro de historia clínica, registro posconsulta, cobro y agendamiento de citas. La comparación de ambas mediciones permite determinar las variaciones registradas en los tiempos de gestión.

La investigación documental se fundamenta en la revisión y análisis sistemático de fuentes bibliográficas, normativas y literatura científica previa. Esta modalidad permite sustentar teóricamente el estudio, establecer antecedentes relevantes y respaldar los criterios de análisis a partir del conocimiento acumulado.

La investigación de campo se centra en la recolección directa de datos primarios en el entorno donde se manifiesta la problemática. Se lleva a cabo mediante la aplicación de entrevistas, encuestas y fichas de observación dirigidas al personal y a los usuarios del establecimiento, capturando la realidad operativa del contexto estudiado.

En conjunto, estos componentes permiten abordar el problema desde una perspectiva integral, respaldar conceptualmente el estudio mediante la evidencia existente y evaluar los cambios observados en el entorno real tras su implementación.